\begin{tabularx}{\textwidth}{|p{0.35\textwidth}|X|}
    \hline \textbf{Notions et contenus} & \textbf{Capacités exigibles} \\
        \hline 1.2. Superposition d'ondes lumineuses & \\
        \hline Superposition de deux ondes quasi-monochromatiques non synchrones ou incohérentes entre elles. & Justifier et utiliser l'additivité des intensités. \\
        \hline Superposition de deux ondes quasi-monochromatiques cohérentes entre elles: formule de Fresnel. & Établir la formule de Fresnel. Identifier une situation de cohérence entre deux ondes et utiliser la formule de Fresnel. \\
        \hline Contraste. & Associer un bon contraste à des ondes d'intensités voisines. \\
        \hline Superposition de $N$ ondes quasi-monochromatiques cohérentes entre elles, de même amplitude et dont les phases sont en progression arithmétique dans le cas $N > 1$. & Expliquer qualitativement l'influence de $N$ sur l'intensité et la finesse des franges brillantes observées. Établir, par le calcul, la condition d'interférences constructives et la demi-largeur $2\pi/N$ des franges brillantes. Établir et utiliser la formule indiquant la direction des maxima d'intensité derrière un réseau de fentes rectilignes parallèles. \\
        \hline
\end{tabularx}
